\chapter{Similarity with the Linux kernel}
\label{cap:similaridade}

\section{What ``similarity'' means}

Comparing this project with Linux requires separating three different questions:

\begin{enumerate}
  \item \textbf{Functional similarity}: which Linux services exist
        here, and with what depth?
  \item \textbf{Architectural similarity}: do the boot, memory, and driver
        solutions use the same concepts and hardware mechanisms?
  \item \textbf{Code similarity}: was there copying, line-by-line
        translation, or reuse of algorithms?
\end{enumerate}

The answers are quite different. The project is conceptually close to Linux
on the boot path and in its use of VGA/PS/2, but it covers a tiny fraction of
the complete kernel. There is no copied code: the implementation is
independent and written in Rust.

\section{Methodology}

The comparison uses as reference the Linux 7.3.1 tree present in
\code{linux-7.3.1/}. The following were collected:

\begin{itemize}
  \item physical lines of the 12 closest Linux files;
  \item physical lines of the 19 Rust modules;
  \item features present, missing, and partial;
  \item hardware mechanisms used (Multiboot2, GDT, PAE, long mode,
        VGA text, 8042, UART 16550);
  \item presence or absence of larger subsystems (IDT, allocator,
        scheduler, syscalls, FS, networking).
\end{itemize}

The score is a grounded estimate, not a formal metric. It serves to convey
order of magnitude and to avoid both the naive conclusion that ``it is Linux
in Rust'' and the opposite conclusion that ``it has nothing to do with it''.

\section{File mapping}

\begin{longtable}{@{}p{7.2cm}r p{4.3cm}@{}}
\caption{Analogous Linux files and corresponding Rust modules.}
\label{tab:correspondencia}\\
\toprule
\textbf{Linux 7.3.1 file} & \textbf{Lines} & \textbf{Rust equivalent} \\
\midrule
\endfirsthead
\toprule
\textbf{Linux 7.3.1 file} & \textbf{Lines} & \textbf{Rust equivalent} \\
\midrule
\endhead
\code{arch/x86/boot/header.S} & 619 & \code{boot.rs} (header) \\
\code{arch/x86/boot/pm.c} & 125 & \code{boot.rs} (stub) \\
\code{arch/x86/kernel/head\_64.S} & 686 & \code{boot.rs} (long mode) \\
\code{arch/x86/kernel/cpu/common.c} & 2,688 & \code{boot.rs} (GDT) + \code{arch.rs} \\
\code{arch/x86/mm/init\_64.c} & 1,557 & \code{boot.rs} (pages) \\
\code{drivers/video/console/vgacon.c} & 1,226 & \code{vga.rs} + \code{tty.rs} \\
\code{drivers/input/serio/i8042.c} & 1,664 & \code{keyboard.rs} \\
\code{drivers/input/keyboard/atkbd.c} & 1,981 & \code{keyboard.rs} \\
\code{drivers/tty/serial/8250/8250\_port.c} & 3,638 & \code{serial.rs} \\
\code{kernel/printk/printk.c} & 5,188 & macros + \code{tty.rs} \\
\code{kernel/panic.c} & 1,274 & \emph{panic handler} in \code{main.rs} \\
\code{init/main.c} & 1,706 & \code{kernel\_main} in \code{main.rs} \\
\code{mm/slab\_common.c}, \code{mm/slub.c} & --- & \code{heap.rs} \\
\code{fs/ramfs/inode.c} & --- & \code{ramfs.rs} \\
\code{fs/open.c}, \code{fs/file\_table.c} & --- & \code{vfs.rs} \\
\code{fs/proc/} & --- & \code{proc.rs} \\
\code{kernel/sched/core.c} & --- & \code{task.rs} \\
\code{arch/x86/entry/syscalls/syscall\_64.tbl} & --- & \code{syscall.rs} \\
\midrule
\textbf{Linux total (first 12 files)} & \textbf{22,352} & --- \\
\textbf{Rust total (19 modules)} & \textbf{7,973} & --- \\
\bottomrule
\end{longtable}

The rows marked \code{---} extend the conceptual mapping after the heap, VFS,
\code{/proc}, tasks and file syscalls were added; they are not part of the
line-count comparison, which still uses the first twelve Linux files.

The line ratio is:

\[
\frac{7{,}973}{22{,}352} \approx 35.67\%.
\]

Relative to the complete Linux tree, which has approximately
38,062,672 lines of \code{.c}, \code{.h}, and \code{.S} in the version used,
the ratio is:

\[
\frac{7{,}973}{38{,}062{,}672} \approx 0.0209\%.
\]

\begin{atencao}[Lines of code do not measure similarity]
The 12 Linux files do not have the same scope as the 19 Rust modules.
\code{vgacon.c}, for example, supports several architectures, virtual
consoles, and graphics modes; \code{vga.rs} implements only 80x25 text.
Moreover, a large part of the Linux lines consists of comments,
configuration, and handling of hardware cases. The LOC comparison serves
only as an order of magnitude.
\end{atencao}

\section{Functional similarity by dimension}

The table below distributes 100 points among ten dimensions and assigns a
score from 0 to 100 for how much of the project covers the corresponding
Linux dimension. The last column is the weighted contribution.

\begin{table}[H]
\centering
\caption{Functional and architectural similarity score.}
\label{tab:score}
\begin{tabular}{@{}lcrrr@{}}
\toprule
\textbf{Dimension} & \textbf{Weight} & \textbf{Linux} & \textbf{Project} & \textbf{Contrib.} \\
\midrule
Boot and firmware protocol & 10 & 100 & 30 & 3.00 \\
Mode transition and GDT & 8 & 100 & 35 & 2.80 \\
Paging and MMU & 12 & 100 & 25 & 3.00 \\
Console and video & 15 & 100 & 25 & 3.75 \\
Keyboard and input & 12 & 100 & 18 & 2.16 \\
Serial and tty & 5 & 100 & 20 & 1.00 \\
Interrupts and exceptions & 12 & 100 & 30 & 3.60 \\
Memory and heap & 8 & 100 & 25 & 2.00 \\
Processes, scheduler, and syscalls & 12 & 100 & 45 & 5.40 \\
Files, networking, and other drivers & 6 & 100 & 10 & 0.60 \\
\midrule
\textbf{Total} & \textbf{100} & --- & --- & \textbf{27.31} \\
\bottomrule
\end{tabular}
\end{table}

The weighted index is \textbf{27.3\%}. It should be read as:

\begin{quote}
``Within the ten chosen dimensions, the project reproduces about 27.3\% of
what Linux offers, with greater weight on boot, console, and keyboard.''
\end{quote}

It does not mean that 27.3\% of the Linux code was reimplemented. In
number of subsystems and services, the fraction is much smaller: the
project has only a first-fit heap, a timer-preemptive weighted fair scheduler
with \code{nice}, weights, \code{vruntime}, a blocking \code{sleep}, event
wait queues, Unix-style signals and zombie reaping, an in-memory filesystem
and a synthetic \code{/proc}, while Linux also has a buddy allocator, slab,
full CFS/EEVDF with real-time classes, a much richer signal and wait queue
infrastructure, a full VFS with block devices, and networking.

\subsection{Reading by dimension}

\begin{itemize}
  \item \textbf{Boot (30/100)}: the project uses Multiboot2 and GRUB, whereas
        Linux uses its own boot protocol, real-time setup, decompression, and
        EFI stubs. The concept of ``load an ELF and transfer control'' is the
        same; the implementation is much simpler.
  \item \textbf{Long mode and GDT (35/100)}: the sequence PAE $\rightarrow$
        LME $\rightarrow$ PG is the same as in \code{head\_64.S}; there is a
        GDT with ring 3 and TSS, but \emph{per-CPU}, \emph{stack canary},
        \emph{paravirt}, and secondary CPU handling are missing.
  \item \textbf{Paging (25/100)}: there are 2 MiB \emph{huge pages} at
        boot, a bitmap frame allocator, 4 KiB tables, and demand
        paging; the buddy allocator, \code{vmalloc}, \code{mmap}, COW,
        swap, NUMA, and \code{NX} protection are missing.
  \item \textbf{Console (25/100)}: VGA text mode and scrolling are
        analogous to \code{vgacon}, but virtual console, fonts,
        Unicode, \emph{fbcon}, \code{printk} with levels, and log buffer are missing.
  \item \textbf{Keyboard (18/100)}: access to the 8042 and scancode
        translation are similar, but IRQ1, \emph{input core},
        \code{evdev}, LEDs, layouts, and extended keys are missing.
  \item \textbf{Serial (20/100)}: the basic programming of the 16550 is
        similar, but the \emph{tty layer}, \emph{termios}, IRQs, input
        FIFO, DMA, and flow control are missing.
  \item \textbf{Interrupts (30/100)}: there is an IDT, \emph{stubs} for
        exceptions, a remapped PIC, and the PIT on IRQ0; APIC, IRQ1,
        spurious IRQs, \emph{softirq}, workqueues, and preemption are missing.
  \item \textbf{Memory/heap (25/100)}: there is a frame bitmap allocator,
        demand paging and a 1 MiB first-fit heap behind
        \code{\#[global\_allocator]}, but there is no buddy allocator, slab,
        page cache, \code{vmalloc}, swap or NUMA.
  \item \textbf{Processes (45/100)}: there is a task table with PID/PPID and
        states, a context switch that preserves \code{RFLAGS}, a
        timer-preemptive weighted fair scheduler with \code{nice},
        \code{sched\_prio\_to\_weight}-style weights and \code{vruntime}, a
        blocking \code{sleep} with \code{schedule\_timeout} semantics, event
        wait queues with \code{wait\_event}/\code{wait\_event\_timeout} and
        \code{wake\_up}, Unix-style signals with pending/blocked masks,
        dispositions and default actions, stopped/continued tasks, exit
        status, zombies, \code{wait}/\code{SIGCHLD} and twelve syscalls;
        there is no \code{fork}/\code{exec}, separate address spaces,
        real-time classes or a full CFS/EEVDF red--black tree.
  \item \textbf{FS/networking (10/100)}: there is a ramfs, a VFS with sixteen
        descriptors and a synthetic \code{/proc}, but there is no block layer,
        disk driver, mount table, permission model, page cache, TCP/IP or
        netfilter.
\end{itemize}

\section{Architectural similarity}

The following table compares concrete mechanisms, regardless of code size.

\begin{table}[H]
\centering
\caption{Comparison of hardware mechanisms and architecture.}
\label{tab:mecanismos}
\begin{tabularx}{\linewidth}{@{}p{3.6cm}p{4.7cm}X@{}}
\toprule
\textbf{Mechanism} & \textbf{Linux 7.3.1} & \textbf{Rust VGA Kernel} \\
\midrule
Loading & Linux protocol + setup + decompression & Multiboot2 via GRUB \\
Initial entry & \code{startup\_32}/\code{startup\_64} & \code{\_start} in assembly \\
Long mode & PAE + LME + \code{CR0.PG} & Same sequence, simplified \\
GDT & \code{gdt\_table}, \emph{per-CPU}, TSS & 7 descriptors, ring 3 and
TSS/IST \\
Paging & 4/5 levels, allocators, \code{vmalloc} & 4 levels,
2 MiB \emph{huge pages} + 4 KiB leaves on demand \\
Console & \code{vgacon}, VT, \code{fbcon}, \code{printk} & \code{vga.rs}
renderer + \code{tty.rs} with eight consoles \\
Keyboard & i8042 + serio + atkbd + input & 8042 by polling (IRQ1 masked) \\
Serial & 8250 + tty + termios & 16550 by polling \\
Input & IRQ1/IRQ4, \emph{softirq}, workqueues & IRQ0 (timer) + polling in the
shell + \code{int 0x80} \\
Time & PIT/APIC, \code{jiffies}, timers & PIT at ~100 Hz, tick
counter, weighted \code{vruntime} accounting \\
Memory & buddy, slab, \code{vmalloc}, page cache & Frame bitmap,
demand paging, first-fit heap \\
Processes & CFS/EEVDF, namespaces, cgroups & Task table, preemptive
weighted fair scheduler (\code{nice}/weight/\code{vruntime}), blocking
\code{sleep}, wait queues, signals and \code{wait}/zombies; no \code{fork} or
namespaces \\
Syscalls & \code{syscall}/\code{sysret}, \code{ptrace} & \code{int 0x80} with
twelve calls (\code{exit}..\code{kill}) \\
Files & VFS, inodes, dentries, page cache & ramfs, 16 descriptors,
synthetic \code{/proc}; no block layer \\
Networking & sockets, TCP/IP, netfilter & Absent \\
\bottomrule
\end{tabularx}
\end{table}

\section{Code similarity}

\begin{itemize}
  \item \textbf{Code copying}: none. The project was written in Rust
        and contains no translated excerpts from Linux C or assembly.
  \item \textbf{Algorithms}: the initialization sequences of the 8042,
        UART, and long mode follow the hardware documentation and consolidated
        patterns, which Linux also follows. This is technical
        convergence, not copying.
  \item \textbf{Constants}: values such as \code{0x3F8}, \code{0x60},
        \code{0x64}, \code{0x3D4}, \code{0xB8000} and the bits of \code{CR0},
        \code{CR4}, and \code{EFER} are defined by the hardware and appear in
        any driver, including Linux.
  \item \textbf{License}: the project declares itself MIT/Apache-2.0. If
        GPL-2.0 excerpts from Linux were copied, the license would need to be
        compatible with GPL. Since the implementation is
        independent, there is no such obligation.
\end{itemize}

\section{Overall similarity}

\begin{table}[H]
\centering
\caption{Summary of the different similarity measures.}
\label{tab:resumo-similaridade}
\begin{tabularx}{\linewidth}{@{}lX@{}}
\toprule
\textbf{Measure} & \textbf{Result} \\
\midrule
Copied code & 0\% (independent implementation) \\
Lines vs.\ 12 analogous files & $\approx 35.67\%$ \\
Lines vs.\ complete Linux tree & $\approx 0.0209\%$ \\
Weighted functional similarity (10 dimensions) & $\approx 27.3\%$ \\
Coverage of complete kernel subsystems & very low; on the order of 2\% \\
Conceptual similarity in boot & high (same x86 mechanisms) \\
Conceptual similarity in drivers & medium/low (polling subset) \\
Conceptual similarity in services & low (basic MM, scheduler, VFS; no net) \\
\bottomrule
\end{tabularx}
\end{table}

\subsection{Classification}

On a qualitative scale:

\begin{itemize}
  \item \textbf{Relationship with Linux}: ``inspired by the same
        concepts and the same hardware'', neither a ``port'' nor a ``clone''.
  \item \textbf{Overall similarity}: \textbf{low}.
  \item \textbf{Educational value}: high for understanding boot, long mode,
        VGA, PS/2, and serial; low for understanding Linux as a complete
        operating system.
  \item \textbf{Score from 0 to 10 for ``how much of Linux was reimplemented''}:
        approximately \textbf{2.7/10} on the weighted score; less than
        \textbf{0.1/10} if the criterion is the complete kernel in
        terms of features.
\end{itemize}

\begin{nota}[How to use this analysis]
If the goal is to study Linux, this project serves as a ``map of the initial
path'': it shows, in a small amount of code, what happens before the Linux
kernel begins to use its subsystems. If the goal is to replace Linux, the
project is several orders of magnitude away: it lacks CFS/EEVDF,
\code{fork}/\code{exec}, per-process address spaces, a disk filesystem,
networking, drivers, and security hardening.
\end{nota}

\section{Essential differences}

\begin{table}[H]
\centering
\caption{Essential differences between the project and Linux.}
\label{tab:diferencas}
\begin{tabularx}{\linewidth}{@{}p{3.2cm}XX@{}}
\toprule
\textbf{Aspect} & \textbf{Linux} & \textbf{Rust VGA Kernel} \\
\midrule
Language & C and assembly & Rust and assembly \\
Size & $\sim$38 million lines & 7,973 lines \\
Boot model & Proprietary protocol, setup, EFI & Multiboot2/GRUB \\
Interrupts & IDT, PIC/APIC, timers & IDT, PIC/PIT; keyboard by polling \\
Concurrency & SMP, preemption, RCU & One CPU, timer-preemptive weighted fair
scheduler with blocking sleep and wait queues \\
Memory & Buddy, slab, VM, page cache & Frame bitmap, demand paging,
first-fit heap \\
Processes & Scheduler, signals, namespaces, cgroups & Task table, preemptive
weighted fair scheduler, wait queues, signals and \code{wait}/zombies; no
namespaces \\
Input/output & VFS, block, net, tty & ramfs, VFS, \code{/proc}, direct I/O ports \\
Security & Rings, SMEP/SMAP, KASLR, LSM & Rings 0/3 and U/S; no NX,
SMEP, or SMAP \\
Portability & Several architectures & x86\_64 BIOS only \\
License & GPL-2.0 & MIT/Apache-2.0 \\
\bottomrule
\end{tabularx}
\end{table}
